This section describes the findings from case B, where the two interviewees from the organization are the IT security manager and the supply chain manager, both working in the IT department. The interviewee from Supplier 1 is a technician in basic IT operations as well as being a security coordinator between Supplier 1 and the organization. The interviewee from Supplier 2 is the service manager for one of Organization B's systems.

\subsection{Preparation}
Organization B defines \textit{security breaches} as events violating the organization's security principles caused intentionally by employees. Examples include employees visiting illegal websites or sharing their passwords. \textit{Security incidents} on the other hand, are events threatening information security, but not necessarily caused by disloyal employees. Security breaches and incidents are often seen as deviations from normal activity and often require some sort of structural change. The organization has a deviation system that keeps track of all kinds of deviations, whether it involves physical shortcomings or issues related to information security.

The interviewee from Supplier 1 defines a security incident related to monitoring of a network as when malware has entered the systems and has the potential to exercise either limited or massive damage. Other security incidents are related to users of the systems, such as disloyal employees, and are much harder to detect. He claimed it is almost impossible. This definition is known among those working with security related to Organization B. The interviewee from Supplier 2 defines unauthorized access of information as being a security incident when this access can cause harm to the organization, customers or other related parties. This is the interviewee's personal definition, but it is grounded in a definition internal to Supplier 2.

When the interviewees from Organization B were asked to consider the worst possible security incident the organization could experience, they said that their biggest fear is disloyal employees leaking information to outsiders without being discovered. Compromised or disclosed information is considered very serious, but in most cases where this is caused by an error it is easier to detect than disloyal employees. Information is the most important asset to protect for the organization as most other things can be restored or repaired.

In addition to information disclosure, they see service interruption as a severe consequence of incidents. Ensuring service availability and business continuity are thus high priorities. Unavailable services could lead to employees not getting any work done which may cause a major financial loss for the organization. Availability is therefore extremely important and thus investments to ensure availability are justified by the financial costs of unavailable services. Redundant equipment and other mechanisms are implemented to ensure the availability of services. Supplier 1 makes sure that Organization B's systems are backed up daily and performs tests to confirm that these backups are usable. They have two data centres in order to increase the level of redundancy.

In addition to financial losses and information leakage, the organization sees reputational damage as a serious consequence of incidents: 

\begin{quote}
\textit{``If sensitive information concerning customers is leaked or lost, it is a contractual breach as well as a violation of trust that would imply risk of reputational damage."}
\end{quote}

The interviewee from Supplier 1 sees loss of sensitive data that are important to the organization's core activities as being the worst possible incident Organization B could experience. He named financial losses, breach of trust and reputational damage as consequences of such incidents. The interviewee from Supplier 2 sees industrial espionage as being the worst security related incident organization B can experience. Such an incident can have financial consequences, as foreign organizations could develop Organization B's products at a lower cost. 

The security policy is the main governing document for activities conducted within Organization B. Its intention is to define senior management's position concerning IT security, and give an overall picture of where the organization stands with regard to information security. They have no policy that specifically addresses incident management, only practical routines and supporting tools.

To avoid vulnerabilities in legacy software being exploited, Organization B's computers are scanned regularly and software with security holes or not in use is removed. When new vulnerabilities are discovered they are categorized and prioritized and Supplier 1 determines how urgent the following update is. For zero-day vulnerabilities they may shut down certain services.

The \acs{IT} manager said that one of the most challenging parts of incident management is constructing a holistic plan. More often than not, things no one had thought of occur and they need to be handled. Hence, ensuring a solid information collection during incidents and making correct decisions at the right time are more important than having a detailed plan to follow. Their experience indicates that the most important thing is to scale correctly, understand the situation and put relevant measures into action. That is also why they believe it is important that not too many have the authority to make important decisions during incidents, which they have tried to limit through their contingency plan.

Organization B has developed plans for communication during incident handling. The contingency plan states who is responsible for communication, both internally and externally. It is always a representative from organization B that is to communicate with the media or the police, not any of the suppliers. The IT manager or the supply chain manager are always available in case of an emergency. To ensure communication during serious incidents, important contact persons have an instant messenger application as well as SIM-cards from several network providers. Organization B has a general contingency plan as well as supplier specific contingency plans and a specific plan for IT services and infrastructure. The IT contingency plan is connected to both the general plan and the supplier specific plans. The IT contingency plan is initiated when the IT manager in consultation with the management and their IT operations suppliers define an incident as a crisis or a catastrophe according to the definitions in the plan.

Supplier 1 keeps track of trends related to security incidents, by monitoring their own internal systems. These trends are discussed with Organization B. Supplier 1 has a case management system where all cases are logged. This way they keep track of previous incidents for Organization B. They also have specific contingency plans developed in cooperation with the organization. These plans are audited every second year. As preventive work, Supplier 1 is responsible for securing the network.

Supplier 2 has a plan for incident management for Organization B as well. The plan is based on various impact levels of incidents. This plan is only related to applications, as Supplier 2 is not responsible for data. They need plans to be in compliance with the \ac{SLA}. They have specific plans regarding communication that specify who to call, when to call, how often to call and what the conversations should contain. Their PR department handles communication with media. Supplier 2 is required to try to fix any vulnerabilities when they are discovered. They are relatively new as a supplier for Organization B and thus have no records of previous incidents.

\paragraph{Standards and Guidelines}
Organization B bases many of their processes on the ITIL framework. Standards are mainly used as a basis for the fundamental thinking related to security in the organization. All IT managers have had training in ISO/IEC 27001 and ISO/IEC 27002 as well as ISO/IEC 27035. Additionally, some employees are certified in ITIL security. Organization B has set compliance with \ac{ITIL} as a requirement for Supplier 2. %Finne ut om dette også er et krav for supplier 1?

Supplier 1 bases everything on \acs{ISO} standards and tries to adapt their contingency plans to these standards. They have not implemented any standards specific to incident management, such as \acs{ISO}/\acs{IEC} 27035. 

The IT manager said that just as important as being familiar with standards is being familiar with the internal documents describing how the organization performs incident management. 

He emphasized that:
\begin{quote}
\textit{``The most important thing when a crisis occurs is knowing what to do, not knowing what the standard says."}
\end{quote}
It is important that employees are familiar with internal routines, hence rehearsals are conducted regularly.

\paragraph{Awareness and Training}
All employees participate in an introductory course where routines for reporting incidents are explained. Additionally, employees are informed through the intranet in cases where they need to be aware of new trends or specific spam e-mails. To raise awareness around IT security, Organization B has previously conducted online lectures for employees via the intranet, addressing various security related topics such as secure use of USB sticks, viruses, social engineering and spam. 

The supply chain manager said that it is difficult to measure the effect of awareness campaigns. Nevertheless, the number of security incidents has decreased during the last three years. 

To raise awareness and best prepare for incident handling in practice Organization B conducts rehearsals regularly, discusses internal routines with employees and includes external suppliers in their training. Both Supplier 1 and Supplier 2 are included in contingency rehearsals. Previously, rehearsals have been set up such that incidents escalate and change as employees discuss what to do in given scenarios. 

IT managers within Organization B as well as the crisis team, basic operations and customer service from Supplier 1 can be involved in rehearsals. They will be presented with an incident and the rehearsal seeks to reveal whether their routines work well or not. It is normally Organization B that initiates these rehearsals. The interviewee from Supplier 1 emphasized the importance of rehearsals:

\begin{quote}
\textit{``Yes, we have to continuously rehearse our routines, otherwise they would never work."}
\end{quote}

Supplier 2 conducts contingency rehearsals on paper for the crisis team. Additionally they have training for the service desk employees. They have as a requirement that employees are to be \ac{ITIL} certified.

Information management, allocation of responsibilities, communication and crisis communication have been identified as areas of improvement after rehearsals. They have revealed the need for a way to communicate to employees during emergency situations that is not dependent on e-mail or the intranet, as these are not necessarily available during a major crisis.

Organization B has outsourced services with several external suppliers and thus collaboration and coordination are extremely important factors in their incident management. Having effective and sufficient coordination as well as distinct and well-established roles during incident response are highlighted by the IT manager as important things to train for. This is supported by the interviewee from Supplier 2 who mentioned that if rehearsals reveal ambiguity concerning roles and responsibilities, they need to change the contracts in order to clarify this.

\subsection{Detection and Analysis}
Organization B's network supplier or other partners may detect incidents. The supply chain manager is convinced that underreporting of incidents exist, but does not see it as a problem. This is supported by the interviewee from Supplier 1. He also mentioned that it is probably impossible to detect everything. The potential high degree of underreporting does not imply that they miss seeing what the main threats are. The supply chain manager said that even though not all incidents are reported, the overall trend is still very apparent in their statistics and thus the right measures to avoid serious incidents can be implemented. 

\paragraph{Initial Detection}
Organization B detects incidents in several ways. Anti-virus detects and reports viruses, and in some cases employees themselves report that their computer is not working as expected. In the latter case the employees are to report to the service desk at Supplier 1. Supplier 1, that has the overall responsibility for monitoring the network, may also detect changes in traffic indicating security incidents. This sensor network is further outsourced to an external security company. The IT manager said that most of these initial detections are handled automatically, and that very few incidents require manual responses. 

\paragraph{Categorization}
Organization B has developed their own framework for categorizing incidents. Inputs to the categorization are the number of persons and departments that are affected in addition to the incident's severity level. The IT manager emphasized that this categorization is quite similar to what is presented in standards and that it is often the availability aspect that is in focus for incident categorization. The central contingency plan states whether incidents should be categorized as a crisis or a catastrophe to ensure that they are handled correctly.

Supplier 2 categorizes an incident to be a crisis if it concerns most of the users and a system is down. This categorization is based on the contract between Supplier 2 and Organization B. The categorization is further used to prioritize incidents.

\subsection{Incident Response}
Minor incidents such as infected computers happen quite often and are seen as part of normal operations. If there is a risk of escalation, users' network ports can be shut down before users are contacted. Sometimes situations arise indicating a serious security breach and the IT manager is contacted. In case of security breaches, the employee involved has to be dealt with which might need involvement from the management.

Organization B does not have an established check-list to follow during incident response. There is a check-list included in the central contingency plan, but not one that addresses IT specifically. This was one of the things that were discovered through a rehearsal and they are currently working on developing such a check-list. The IT manager said this is an area they wish to further develop to improve their incident handling capabilities. 

Lack of proper communication may lead to unnecessary trouble. The supply chain manager said that the implicit knowledge of responsibilities in minor cases is an example of routines that are difficult to document properly.

\begin{quote}
\textit{``Problems may arise even with minor incidents when everyone assumes that it is someone else's responsibility."}
\end{quote}

The supply chain manager highlighted the challenge of deciding in what way and how much information should be given to employees during incident handling. This is challenging as it is very individual how much information people both want to share and receive, he said. Further, the IT manager believes knowing when to communicate, what to say and to whom might be the most challenging aspects of incident management.

When Supplier 1 first detects incidents they usually notify the organization by e-mail. After detection they initialise a team that handles the incident. If the incident spreads there can be assembled a crisis team and the communication in that case is by e-mail, phone and personal contact. Most cases that Supplier 1 handles are routine cases and standard procedures can be used. If an incident is very serious the security coordinator is involved and is responsible for communication between Supplier 1 and the organization. All incidents are logged in the same system. The log includes root cause of the incident, the impact of it, potential breaches of \acp{SLA}, actions taken and solutions to the incident. 

Supplier 2 has procedures for handling known security incidents. They do not have any automatic processes as incidents suited for automatic handling are handled by Supplier 1. Many of the activities in incident handling are determined as the incident evolves. This is based on what is experienced at that time, previous experiences, feedback and expectations. To ensure that incidents are solved a fast as possible, the technician(s) who gets the task will only focus on fixing the incident and is not to be disturbed until it is solved. 

\paragraph{\acl{IRT}}
Organization B does not have a dedicated IRT, but they have a team responding to incidents that is dynamic, it changes with the type of incident they are dealing with. Supplier 1 acts as a response team and handles most minor incidents related to IT operations as part of their normal activities. The responsible person and the response team depend on the characteristics of the incident.   

For minor incidents, the team is almost entirely outsourced to Supplier 1, but is gradually insourced as the severity of the incident increases. Thus, members of the team may vary. The IT manager said:

\begin{quote}
\textit{``We try to scale the team in response to the specific incident we are dealing with."}
\end{quote}

The sequence of the scaling is stated in the IT contingency plan, where particular roles and activities are described. The permanent members of the team are the IT manager and the supply chain manager as they are involved in the handling of most incidents. Team members have other tasks beside incident management, but if there is need for them to aid in incident response, all other tasks are put on hold.

For incidents categorized as crises, a crisis team will be formed. The team is composed of the IT manager, the supply chain manager, management from Supplier 1, the person at Supplier 1 responsible for IT operations for Organization B at the time of the crisis and whenever necessary representatives from Supplier 2. For the most serious incidents, a central manager in Organization B will be in charge of the situation.

Supplier 1 has a team that is available 24/7 and receives (often automated) incidents reports. They can subsequently determine what to do about it. Supplier 2 also has teams to handle incidents for Organization B. They have a crisis team and a support function to handle daily management and maintenance. The members of the team are available 24/7. They look for improvements in the process and have ongoing reviews of plans. 

\paragraph{Workflows}
%The organization uses a proactive method for collecting information during incidents. %beskrevet bakerst i den sentrale beredskapsplanen.

Supplier 1 handles incidents related to IT operations. Figure \ref{fig:WorkflowCaseBSupplier1} illustrates the workflow for incidents (that are not escalated and categorized as crises).

\begin{figure}[H]
%\hspace{-1.1cm}
\begin{center}
\includegraphics[scale=0.57]{WorkflowCaseBSupplier1.png}
\caption[Workflow for Incidents, Case B Supplier 1]{Workflow for Incidents for Supplier 1}
\label{fig:WorkflowCaseBSupplier1}
\end{center}
\end{figure}

\begin{itemize}
\item The external supplier of the sensor-network receives notifications from sensors and categorizes the incidents.
\item If the incident is categorized as high the customer service centre at Supplier 1 gets a notification as well as a phone call from the external supplier of the sensor-network. The customer service centre can also receive notifications from users.
\item The customer service centre investigates each individual case.
\item The network access for the equipment in question is removed 
\item The user and the IT manager of the department involved are notified.
\item The equipment is brought in to Supplier 1 for inspection and subsequently fixed.
\end{itemize}

There are dedicated people that take care of this process and there are 2-3 people on rotation. These people are well informed of the procedure.  

\begin{figure}[H]
%\hspace{-1.1cm}
\begin{center}
\includegraphics[scale=0.57]{WorkflowCaseBSupplier2.png}
\caption[Workflow for Incidents, Case B Supplier 2]{Workflow for Incidents for Supplier 2}
\label{fig:WorkflowCaseBSupplier2}
\end{center}
\end{figure}

Supplier 2 is responsible for handling incidents related to applications in Organization B. Their incident workflow is illustrated in figure \ref{fig:WorkflowCaseBSupplier2}.

\begin{itemize}
\item The service desk receives incident reports. Supplier 2 is the second line of incident handling, so they do not get reports directly from users. Usually they receive reports from Supplier 1.
\item The severity of the incident is assessed by the service desk employee in cooperation with the person reporting.
\item If the incident is critical the service desk calls a person responsible for escalation who subsequently initiates the crisis management team. This person can call people and have them come to work if necessary. The crisis team is responsible for getting the incident solved.
\item For any other incident the service desk will document all reported information and transfer the case to the correct group of technicians.
\end{itemize}

\paragraph{Escalation}
Organization B's team is scalable, but it escalates whenever incidents escalate and reach a high severity level. Whenever incidents are of such severity that escalation is necessary, they primarily try to recruit extra internal employees to the team. In some cases external experts are needed to respond effectively to an incident.

For incidents that are related to employee privacy, Supplier 1 calls in a central security advisor. This is done to avoid privacy violations. Supplier 2 has developed clear routines for escalation and contact persons both in Organization B and Supplier 1. They have mandate to call in an external crisis management team if necessary. 

\paragraph{Electronic Evidence}
In case of security breaches, i.e. employees violating policies deliberately, all logs are preserved in case of need for future investigations. In cases where it is suspected that an employee is being disloyal, his user account is blocked to avoid potential evidence being deleted. The supply chain manager said assessing what to do in each individual case is difficult as sometimes incidents may be more serious than first assumed. Supplier 1 brings in an expert from central parts of their organization in cases where electronic evidence must be preserved. In these cases all communication must be encrypted. The police can be brought in or evidence can be handed over to the police for further investigations.

\subsection{Lessons Learned}
Incident reports are in most cases constructed by the two main suppliers and include excerpts from logs, information on what happened, what was the cause, what was done to solve the incident and proposed measures to avoid similar incidents in the future. Organization B receives monthly reports from both their main IT suppliers, where all incidents are recorded. Supplier 2 only includes critical incidents or specially requested incidents in their reports. All IT managers as well as the security coordinator from Supplier 1 participate in a monthly meeting where monthly incident reports from the organization's main suppliers are discussed. 

Organization B only conducts debriefing meetings after serious incidents, which might be a couple of times a year. The participants in these meetings are the involved IT manager(s), the service manager and the security coordinator from Supplier 1. Additionally, four times a year, incidents concerning intelligence, espionage, crime for profit and information security are discussed within the organization's security board and are reported to the management. 

The IT manager said that they are interested in obtaining a high contingency and security in their operations and that learning from others is an important factor in achieving that. To benefit from others' experiences they often think about how incidents in other organizations would have affected themselves. In that matter, incidents that occurred in other organizations contribute to development and improvement of contingency plans. The IT manager said that learning from other's experience is important since the scenarios are real and could most likely happen in their organization as well.

The interviewee from Supplier 1 said that when they experience botnets they are able to handle them fast and effectively. He mentioned that one thing they learned from rehearsals is not to trust electronics. Thus, contingency plans should be available on paper as well as electronically. Supplier 2 has so far been able to handle all critical incidents the last years in accordance with the \ac{SLA}. One contributing factor to this success was that they contacted the right technical resources when needed. Nevertheless, routines have not always worked well, and an example is when different suppliers are involved and no one takes responsibility for the error. Supplier 2 is a consortium and therefore they often have to cooperate. The interviewee has experienced that various suppliers have different focuses:

\begin{quote}
\textit{``Some are for example of the understanding that they do not own the error, while someone else understands that they own it. In cases where the owner does not understand it, it is very challenging to make it work. And it takes a lot of time [...] This is the greatest challenge [...] It is also a political ``game". Who will pay for it?"}
\end{quote}

Supplier 2 shares experiences from incident handling internally in their team. They conduct post-review sessions where they discuss what was done, what was done in a bad way and what can be improved. This is done in cooperation with Organization B, but not with other suppliers. The interviewee mentioned that there are many sources of information and handling these and collecting the correct information can be very difficult.

\subsection{Employee Survey}
Among the 15 participating employees, only two answered that they were not familiar with the organization's information security policy. Eight answered that they were to some extent familiar with it, but that they did not know the details or that they had read it once. One of the employees mentioned that there have been made many changes in the policy during the last years and that employees have not always been notified about the changes.

Only one employee had never received any suspicious e-mails. Several said that they receive such e-mails very often, but that it is mainly ``common" spam and easy to recognize. No one mentioned to have received targeted e-mails. No one acknowledged to have carried out instructions in such e-mails, except one that claimed he once did it in a secure way out of curiosity. He clicked a link to examine the quality of the phishing-site, but did not carry out any instructions from there on. A few of the employees have reported to have received such e-mails, but several mentioned that they do not report it, because it is so common. One mentioned that it happens so often, so it is difficult to know when you are suppose to report, but that he would report whenever in doubt.

Six employees claimed they knew what an information security incident is. Five could not provide a definition, but had an idea of what it is. The rest claimed that they did not know. Most of them were however able to give examples. Among them were information or passwords astray, storage of confidential information on unencrypted \acsp{USB} sticks, use of unknown \acsp{USB} sticks, breach of security instructions or sharing your computer. There were also given examples that showed that there were employees who obviously did not know what an information security incident is:

\begin{quote}
\textit{``Information security incident [...] That must be when someone gives away money they are not supposed to?"}
\end{quote}

Most employees claimed they were attentive to incidents, even though several of them stated that they did not know what an incident is. Some elaborated their statement by providing examples. They said that lock their screens when they leave their computers and lock their doors when they leave their offices and destroy sensitive documents. The IT manager of one of the departments said that they try to make sure that employees know that they will not be held responsible if they caused a security incident. This is done to mitigate underreporting of incidents.

Most of the employees were unsure about which cases they should report. Most of them claimed that they would have reported to the local or central IT manager, but some mentioned that they did not know if this was correct. One employee said that he would report to the local IT manager or to Supplier 1. Some of the employees that were not able to define a security incident, still believed they would be able to know whether to report should an incident occur. There were also a few claiming that they were not familiar with reporting procedures, because they had never needed to be.

Organization B conducted an awareness campaign where all employees were to go through a set of online slides each week for a period of time. The IT manager of one of the departments claimed they got positive feedback from employees. This is supported by several of the employees in the survey who acknowledged to have read these slides and thought it was useful, even though much of it was well-known material. Several of the employees want such lectures more often:

\begin{quote}
\textit{``The content was known, but it is all right to have read it."}
\end{quote} 

\begin{quote}
\textit{``Having an IT instruction is one thing, but being \textbf{reminded} of security instructions is always useful."}
\end{quote}

Only two employees who had read the slides in the awareness campaign did not find them useful. One of them mentioned that some of the proposed measures were socially unacceptable, such as refusing a customer to use their own USB sticks when using one of the organization's computers during a presentation. The other employee mentioned that some of the measures were so strict that he would rather take the risk of security incidents than following them.